\documentclass{amsart}
% For dealing with references we use the comment environment
\usepackage{verbatim}
\newenvironment{reference}{\comment}{\endcomment}
%\newenvironment{reference}{}{}
\newenvironment{slogan}{\comment}{\endcomment}
\newenvironment{history}{\comment}{\endcomment}

% For commutative diagrams we use Xy-pic
\usepackage[all]{xy}

% We use 2cell for 2-commutative diagrams.
\xyoption{2cell}
\UseAllTwocells

% We use multicol for the list of chapters between chapters
\usepackage{multicol}

% This is generally recommended for better output
\usepackage{lmodern}
\usepackage[T1]{fontenc}

% For cross-file-references

% Package for hypertext links:
\usepackage{hyperref}

% For any local file, say "hello.tex" you want to link to please

% Theorem environments.
%
\theoremstyle{plain}
\newtheorem{theorem}[subsection]{Theorem}
\newtheorem{proposition}[subsection]{Proposition}
\newtheorem{lemma}[subsection]{Lemma}

\theoremstyle{definition}
\newtheorem{definition}[subsection]{Definition}
\newtheorem{example}[subsection]{Example}
\newtheorem{exercise}[subsection]{Exercise}
\newtheorem{situation}[subsection]{Situation}

\theoremstyle{remark}
\newtheorem{remark}[subsection]{Remark}
\newtheorem{remarks}[subsection]{Remarks}

\numberwithin{equation}{subsection}

% Macros
%
\def\lim{\mathop{\mathrm{lim}}\nolimits}
\def\colim{\mathop{\mathrm{colim}}\nolimits}
\def\Spec{\mathop{\mathrm{Spec}}}
\def\Hom{\mathop{\mathrm{Hom}}\nolimits}
\def\Ext{\mathop{\mathrm{Ext}}\nolimits}
\def\SheafHom{\mathop{\mathcal{H}\!\mathit{om}}\nolimits}
\def\SheafExt{\mathop{\mathcal{E}\!\mathit{xt}}\nolimits}
\def\Sch{\mathit{Sch}}
\def\Mor{\mathop{\mathrm{Mor}}\nolimits}
\def\Ob{\mathop{\mathrm{Ob}}\nolimits}
\def\Sh{\mathop{\mathit{Sh}}\nolimits}
\def\NL{\mathop{N\!L}\nolimits}
\def\CH{\mathop{\mathrm{CH}}\nolimits}
\def\proetale{{pro\text{-}\acute{e}tale}}
\def\etale{{\acute{e}tale}}
\def\QCoh{\mathit{QCoh}}
\def\Ker{\mathop{\mathrm{Ker}}}
\def\Im{\mathop{\mathrm{Im}}}
\def\Coker{\mathop{\mathrm{Coker}}}
\def\Coim{\mathop{\mathrm{Coim}}}

% Boxtimes
%
\DeclareMathSymbol{\boxtimes}{\mathbin}{AMSa}{"02}

%
% Macros for moduli stacks/spaces
%
\def\QCohstack{\mathcal{QC}\!\mathit{oh}}
\def\Cohstack{\mathcal{C}\!\mathit{oh}}
\def\Spacesstack{\mathcal{S}\!\mathit{paces}}
\def\Quotfunctor{\mathrm{Quot}}
\def\Hilbfunctor{\mathrm{Hilb}}
\def\Curvesstack{\mathcal{C}\!\mathit{urves}}
\def\Polarizedstack{\mathcal{P}\!\mathit{olarized}}
\def\Complexesstack{\mathcal{C}\!\mathit{omplexes}}
% \Pic is the operator that assigns to X its picard group, usage \Pic(X)
% \Picardstack_{X/B} denotes the Picard stack of X over B
% \Picardfunctor_{X/B} denotes the Picard functor of X over B
\def\Pic{\mathop{\mathrm{Pic}}\nolimits}
\def\Picardstack{\mathcal{P}\!\mathit{ic}}
\def\Picardfunctor{\mathrm{Pic}}
\def\Deformationcategory{\mathcal{D}\!\mathit{ef}}

\setlength{\emergencystretch}{3em}
\begin{document}
\title[Stacks Project, Tag 0573]{469 --- Scheme-theoretic density in a generic fibre spreads over an open subset of an irreducible base}
\maketitle
\noindent Copyright (C) 2005--2025 Johan de Jong. Stacks Project authors. Source: \href{https://stacks.math.columbia.edu/tag/0573}{Tag 0573}.

\noindent Editorial difficulty note (not author text). Difficulty H1: A topologically dense open need not detect embedded or nilpotent structure, so a finite algebraic injectivity test must be built. The proof encodes the generic density by multiplication with finitely many functions, applies generic flatness to the algebra and cokernel, and transports the exact injection to every fibre on the chosen open..

\noindent Editorial provenance note (not author text). History: excerpt from revision \texttt{a0444\allowbreak 6e57e\allowbreak c1fbc\allowbreak 252a8\allowbreak 71afc\allowbreak ec775\allowbreak 2fb28\allowbreak 07b14}, more-morphisms.tex, lines 6562--6641. Reference presentation was adapted; original-author context and core support blocks were retained or added. The mathematical wording is unchanged. Licensed under GNU FDL 1.2 or later; no invariant sections or cover texts. See ../licenses/COPYING. Context blocks reproduce author definitions and supporting results; they are not additional counted problems.

\begin{definition}
\label{morphisms-definition-scheme-theoretically-dense}
Let $X$ be a scheme. Let $U \subset X$ be an open subscheme.
\begin{enumerate}
\item The scheme theoretic image of the morphism $U \to X$
is called the {\it scheme theoretic closure of $U$ in $X$}.
\item We say $U$ is {\it scheme theoretically dense in $X$}
if for every open $V \subset X$ the scheme theoretic closure
of $U \cap V$ in $V$ is equal to $V$.
\end{enumerate}
\end{definition}

\begin{lemma}
\label{lemma-scheme-theoretically-dense-generic-fibre}
Let $f : X \to Y$ be a finite type morphism of schemes.
Assume $Y$ irreducible with generic point $\eta$.
Let $U \subset X$ be an open subscheme such that $U_\eta$ is
scheme theoretically dense in $X_\eta$.
Then there exists a nonempty open $V \subset Y$ such
that $U_y$ is scheme theoretically dense in $X_y$ for all $y \in V$.
\end{lemma}

\begin{proof}
Let $Y' \subset Y$ be the reduction of $Y$.
Let $X' = Y' \times_Y X$ and $U' = Y' \times_Y U$.
As $Y' \to Y$ induces a bijection on points, and as
$U' \to U$ and $X' \to X$ induce isomorphisms of scheme theoretic fibres,
we may replace $Y$ by $Y'$ and $X$ by $X'$.
Thus we may assume that $Y$ is integral, see
Properties, Lemma \href{https://stacks.math.columbia.edu/tag/01ON}{lemma characterize integral (Tag 01ON)}.
We may also replace $Y$ by a nonempty affine open. In other words we
may assume that $Y = \Spec(A)$ where $A$ is a domain with fraction
field $K$.

\medskip\noindent
As $f$ is of finite type we see that $X$ is quasi-compact.
Write $X = X_1 \cup \ldots \cup X_n$ for some affine opens $X_i$. By
Morphisms, Definition \ref{morphisms-definition-scheme-theoretically-dense}
we see that $U_i = X_i \cap U$ is an open subscheme of $X_i$ such that
$U_{i, \eta}$ is scheme theoretically dense in $X_{i, \eta}$.
Thus it suffices to prove the result for the pairs $(X_i, U_i)$,
in other words we may assume that $X$ is affine.

\medskip\noindent
Write $X = \Spec(B)$. Note that $B_K$ is Noetherian as it is a
finite type $K$-algebra. Hence $U_\eta$ is quasi-compact. Thus we can
find finitely many $g_1, \ldots, g_m \in B$ such that $D(g_j) \subset U$
and such that $U_\eta = D(g_1)_\eta \cup \ldots \cup D(g_m)_\eta$.
The fact that $U_\eta$ is scheme theoretically dense in
$X_\eta$ means that $B_K \to \bigoplus_j (B_K)_{g_j}$
is injective, see
Morphisms, Example \href{https://stacks.math.columbia.edu/tag/056C}{example scheme theoretic closure (Tag 056C)}.
By
Algebra, Lemma \href{https://stacks.math.columbia.edu/tag/0565}{lemma when injective covering (Tag 0565)}
this is equivalent to the injectivity of
$B_K \to \bigoplus\nolimits_{j = 1, \ldots, m} B_K$,
$b \mapsto (g_1b, \ldots, g_mb)$. Let $M$ be the cokernel of this
map over $A$, i.e., such that we have an exact sequence
$$
0 \to I \to B \xrightarrow{(g_1, \ldots, g_m)}
\bigoplus\nolimits_{j = 1, \ldots, m} B \to M \to 0
$$
After replacing $A$ by $A_h$ for some nonzero $h$ we may assume that $B$
is a flat, finitely presented $A$-algebra, and that $M$
is flat over $A$, see
Algebra, Lemma \href{https://stacks.math.columbia.edu/tag/051T}{lemma generic flatness (Tag 051T)}.
The flatness of $B$ over $A$ implies that $B$ is torsion free as an
$A$-module, see
More on Algebra, Lemma \href{https://stacks.math.columbia.edu/tag/0538}{lemma flat torsion free (Tag 0538)}.
Hence $B \subset B_K$. By assumption $I_K = 0$ which implies that $I = 0$
(as $I \subset B \subset B_K$ is a subset of $I_K$). Hence now
we have a short exact sequence
$$
0 \to B \xrightarrow{(g_1, \ldots, g_m)}
\bigoplus\nolimits_{j = 1, \ldots, m} B \to M \to 0
$$
with $M$ flat over $A$. Hence for every homomorphism $A \to \kappa$ where
$\kappa$ is a field, we obtain a short exact sequence
$$
0 \to B \otimes_A \kappa \xrightarrow{(g_1 \otimes 1, \ldots, g_m \otimes 1)}
\bigoplus\nolimits_{j = 1, \ldots, m} B \otimes_A \kappa \to
M \otimes_A \kappa \to 0
$$
see
Algebra, Lemma \href{https://stacks.math.columbia.edu/tag/00HL}{lemma flat tor zero (Tag 00HL)}.
Reversing the arguments above
this means that $\bigcup D(g_j \otimes 1)$ is scheme
theoretically dense in $\Spec(B \otimes_A \kappa)$.
As $\bigcup D(g_j \otimes 1) = \bigcup D(g_j)_\kappa \subset U_\kappa$
we obtain that $U_\kappa$ is scheme theoretically dense in $X_\kappa$
which is what we wanted to prove.
\end{proof}
\end{document}
